\documentclass[11pt]{article}

\usepackage[margin=1in]{geometry}
\usepackage{booktabs}
\usepackage{siunitx}
\usepackage{amsmath}

\sisetup{detect-all=true}

\title{Benchmark Summary for \texttt{aerosolDynamicsFoam}}
\author{}
\date{}

\begin{document}
\maketitle

\section{Purpose}

The benchmarks were performed with the \texttt{testImpactor} case to evaluate
three implementation choices:

\begin{enumerate}
    \item whether \texttt{analyticPositionTracking} permits larger particle
    time steps without losing the relevant particle fate information;
    \item whether the replicated-mesh particle-parallel mode scales well for a
    fixed carrier-flow, one-way Lagrangian calculation;
    \item how the event-driven and ordinary time-loop modes compare.
\end{enumerate}

The main result is that analytic position tracking makes coarse particle time
steps considerably more usable.  In the impaction case, it reduces the
coarse-step bias in the outlet/substrate split, so the calculation can use a
larger \texttt{deltaTMax} while preserving the practically important particle
fate trends.  Full convergence of the fate counts was still obtained only when
the particle step was reduced to about \(\Delta t_{\max}=10^{-7}\,\mathrm{s}\)
for the present case, but the analytic displacement update gives a much better
large-step approximation than the legacy position update.

\section{Benchmark Set-up}

All benchmark cases were derived from \texttt{testImpactor}.  The original case
was copied to a temporary benchmark directory, and the carrier-flow field and
mesh were kept fixed.  The particle size range was
\(0.10\)--\(1.00\,\si{\micro\metre}\), with a spacing of
\(0.05\,\si{\micro\metre}\).  The injection positions were duplicated from the
original valid injection set, giving 200 parcels per size bin and 3800 parcels
in total.

The event-driven calculations used

\[
    t_\mathrm{track}=0.009119\,\mathrm{s}.
\]

The particle-parallel calculations were run with four MPI ranks using
\texttt{-particleParallel}.  This is not OpenFOAM mesh-decomposition parallelism:
each rank reads the full mesh and carrier field, while the injected parcels are
split into particle shards.

\section{Effect of \texttt{analyticPositionTracking}}

Table~\ref{tab:true-event} shows the event-driven results with
\texttt{analyticPositionTracking true}.  The finest case,
\(\Delta t_{\max}=10^{-8}\,\mathrm{s}\), is used here as the reference.

\begin{table}[h]
\centering
\caption{Event-driven benchmark with \texttt{analyticPositionTracking true}.}
\label{tab:true-event}
\begin{tabular}{@{}lrrrrrr@{}}
\toprule
\(\Delta t_{\max}\) & Wall time [s] & outlet & substrate & firstConv & firstNozzle & secondConv \\
\midrule
\(10^{-4}\) & 7.68   & 1266 & 2420 & 76 & 38 & 0  \\
\(10^{-5}\) & 7.65   & 1266 & 2420 & 76 & 38 & 0  \\
\(10^{-6}\) & 7.36   & 1260 & 2426 & 76 & 6  & 16 \\
\(10^{-7}\) & 20.06  & 1316 & 2370 & 76 & 8  & 30 \\
\(10^{-8}\) & 189.82 & 1290 & 2396 & 76 & 8  & 30 \\
\bottomrule
\end{tabular}
\end{table}

At coarse steps, the largest fate-distribution differences relative to the
\(10^{-8}\,\mathrm{s}\) reference occurred on the small-particle side.  Using
the total-variation difference of the patch fate distribution per size bin, the
affected size bins were:

\begin{itemize}
    \item \(10^{-4}\) and \(10^{-5}\,\mathrm{s}\): mainly
    \(0.10\)--\(0.25\,\si{\micro\metre}\), with a maximum difference of about
    8\%;
    \item \(10^{-6}\,\mathrm{s}\): mainly \(0.10\)--\(0.25\,\si{\micro\metre}\),
    with an additional sensitivity around \(0.85\,\si{\micro\metre}\);
    \item \(10^{-7}\,\mathrm{s}\): no size bin exceeded a 5\% difference, and
    the maximum difference was about 2\%.
\end{itemize}

Table~\ref{tab:false-event} shows the same event-driven benchmark with
\texttt{analyticPositionTracking false}.  The same
\texttt{analyticPositionTracking true}, \(10^{-8}\,\mathrm{s}\) case is used as
the reference for the error statements.

\begin{table}[h]
\centering
\caption{Event-driven benchmark with \texttt{analyticPositionTracking false}.}
\label{tab:false-event}
\begin{tabular}{@{}lrrrrrrr@{}}
\toprule
\(\Delta t_{\max}\) & Wall time [s] & outlet & substrate & firstConv & firstNozzle & secondConv & secondNozzle \\
\midrule
\(10^{-4}\) & 1.23 & 1184 & 2502 & 82 & 32 & 0  & 0  \\
\(10^{-5}\) & 1.24 & 1184 & 2502 & 82 & 32 & 0  & 0  \\
\(10^{-6}\) & 1.51 & 1180 & 2506 & 76 & 8  & 20 & 10 \\
\(10^{-7}\) & 9.18 & 1288 & 2398 & 76 & 8  & 30 & 0  \\
\bottomrule
\end{tabular}
\end{table}

Without analytic position tracking, the coarse-step solution over-predicts
substrate collection and under-predicts outlet escape.  For example, at
\(\Delta t_{\max}=10^{-4}\,\mathrm{s}\), the analytic update gives
1266 outlet parcels, whereas the legacy position update gives only 1184 outlet
parcels; the \(10^{-8}\,\mathrm{s}\) reference gives 1290.  Thus the analytic
displacement update substantially reduces the coarse-step bias in the primary
outlet/substrate split.

The \texttt{false} runs also converge when the time step is reduced.  At
\(10^{-7}\,\mathrm{s}\), the total fate counts are very close to the reference:
1288 outlet parcels and 2398 substrate parcels, compared with 1290 and 2396 in
the reference.  However, the benefit of \texttt{analyticPositionTracking} is
that useful accuracy is retained at a larger particle time step.

\section{Comparison with \texttt{icoUncoupledKinematicParcelFoam}}

As a validation check, the same impactor benchmark was also run with the
standard OpenFOAM \texttt{icoUncoupledKinematicParcelFoam} solver.  This solver
does not contain the new event-driven tracking, analytic displacement update, or
particle-sharding options.  It is therefore a useful baseline for checking that
the present implementation recovers the legacy behavior when
\texttt{analyticPositionTracking} is disabled.

Table~\ref{tab:ico-compare} compares the custom solver with
\texttt{analyticPositionTracking false} against
\texttt{icoUncoupledKinematicParcelFoam}.  The fate tuple is
\[
    (N_\mathrm{outlet}, N_\mathrm{substrate}, N_\mathrm{firstConv},
     N_\mathrm{firstNozzle}, N_\mathrm{secondConv}).
\]

\begin{table}[h]
\centering
\caption{Legacy-behavior validation against \texttt{icoUncoupledKinematicParcelFoam}.}
\label{tab:ico-compare}
\begin{tabular}{@{}lllrr@{}}
\toprule
\(\Delta t\) & Custom solver, analytic=false & \texttt{icoUncoupled} & Custom time [s] & \texttt{ico} time [s] \\
\midrule
\(10^{-4}\) & (1184, 2502, 82, 32, 0) & (1184, 2502, 82, 32, 0) & 1.23 & 3.82  \\
\(10^{-5}\) & (1184, 2502, 82, 32, 0) & (1186, 2500, 80, 34, 0) & 1.24 & 4.36  \\
\(10^{-6}\) & (1180, 2506, 76, 8, 20) & (1208, 2478, 76, 8, 30) & 1.51 & 10.43 \\
\(10^{-7}\) & (1288, 2398, 76, 8, 30) & (1288, 2398, 76, 8, 30) & 9.18 & 96.47 \\
\bottomrule
\end{tabular}
\end{table}

The agreement is strongest at the finest tested step:
\(\Delta t=10^{-7}\,\mathrm{s}\) gives identical patch fate counts between the
custom solver with \texttt{analyticPositionTracking false} and
\texttt{icoUncoupledKinematicParcelFoam}.  At coarser steps, the two solvers
show the same qualitative trend, with small differences that decrease as the
time step is refined.  This is an important regression test: disabling the new
analytic displacement update recovers the established solver behavior.  The
remaining improvement with \texttt{analyticPositionTracking true} can therefore
be attributed to the new position update rather than to an unrelated change in
the particle model or impactor set-up.

\section{Event-driven Versus Ordinary Time Loop}

The ordinary non-event-driven time-loop calculations were also run with four
particle shards.  In this mode the carrier time step controls the number of
outer time-loop iterations.  The measured results are shown in
Table~\ref{tab:loop}.

\begin{table}[h]
\centering
\caption{Ordinary time-loop benchmark, without event-driven tracking.}
\label{tab:loop}
\begin{tabular}{@{}lrrrrrr@{}}
\toprule
\(\Delta t\) & Wall time [s] & outlet & substrate & firstConv & firstNozzle & secondConv \\
\midrule
\(10^{-4}\) & 7.63  & 1266 & 2420 & 76 & 38 & 0  \\
\(10^{-5}\) & 7.67  & 1268 & 2418 & 76 & 38 & 0  \\
\(10^{-6}\) & 6.48  & 1290 & 2396 & 76 & 6  & 32 \\
\(10^{-7}\) & 55.70 & 1308 & 2378 & 76 & 8  & 30 \\
\bottomrule
\end{tabular}
\end{table}

The absolute timings include MPI start-up and I/O/logging overhead, so they
should not be over-interpreted as a smooth time-step scaling curve.  The
important algorithmic point is that the ordinary time loop requires many more
outer loop iterations as \(\Delta t\) decreases.  In contrast, event-driven
tracking advances the particles over a prescribed physical tracking horizon
without repeatedly cycling the fixed Eulerian field.

\section{Particle Parallelism}

The solver uses a replicated-mesh particle-sharding strategy.  Each rank reads
the full mesh and carrier fields, and only the particles are divided among
ranks.  This is well matched to the present problem because the carrier flow is
fixed and the calculation is one-way coupled.

A two-shard benchmark on the impactor case gave the following timings:

\begin{table}[h]
\centering
\caption{Two-shard particle-parallel timing check.}
\label{tab:parallel}
\begin{tabular}{@{}lrr@{}}
\toprule
Run & Wall time [s] & Relative work \\
\midrule
Serial & 7.45 & all parcels \\
Shard 0 & 3.73 & half the parcels \\
Shard 1 & 3.75 & half the parcels \\
\bottomrule
\end{tabular}
\end{table}

If the two shards are executed concurrently, the expected speed-up is roughly
\[
    \frac{7.45}{3.75} \approx 1.99,
\]
which is essentially ideal two-way scaling for this particle-dominated
benchmark.  A small-parcel-count MPI test showed much less benefit because the
fixed costs of MPI launch, mesh reading, and output become comparable to the
particle work.

\begin{table}[h]
\centering
\caption{Standard OpenFOAM cell-decomposition timing with \texttt{icoUncoupledKinematicParcelFoam}.}
\label{tab:ico-parallel}
\begin{tabular}{@{}lrr@{}}
\toprule
\(\Delta t\) & Serial time [s] & Four-way mesh-decomposed time [s] \\
\midrule
\(10^{-4}\) & 3.82  & 4.01  \\
\(10^{-5}\) & 4.36  & 4.73  \\
\(10^{-6}\) & 10.43 & 10.49 \\
\(10^{-7}\) & 96.47 & 97.24 \\
\bottomrule
\end{tabular}
\end{table}

Table~\ref{tab:ico-parallel} shows that standard cell decomposition did not
accelerate the baseline \texttt{icoUncoupledKinematicParcelFoam} benchmark.
For this fixed-flow particle-tracking problem, the mesh-decomposition overhead
and particle migration costs do not pay for themselves.  This contrasts with
particle sharding, which directly divides the dominant particle work.

\section{Why Particle Sharding Is Preferable Here}

Standard cell decomposition is generally not the best first choice for this
solver.  The carrier field is fixed, and the expensive part is tracking many
independent particles.  Decomposing cells introduces processor patches and
particle migration between subdomains, but it does not remove the main
particle-tracking cost.  It can also create load imbalance when many particles
collect or escape in some parts of the geometry but not others.

Particle sharding avoids those costs.  Since each rank has the full mesh,
particles do not need to communicate when crossing a processor boundary.  For
one-way tracking on a fixed flow field, this is a natural decomposition of the
work.

\section{Caveats}

Both approaches have caveats.

\paragraph{Replicated-mesh particle sharding.}
Memory usage grows approximately with the number of ranks because every rank
holds the full mesh and carrier fields.  The output is also sharded, so
post-processing must combine \texttt{kinematicCloud\_shard*} results.  The
method is most effective when the parcel count is large enough that particle
tracking dominates the fixed cost of reading the case and writing output.
Load balance is based on the initial particle shard assignment, so cases with
strongly size-dependent or position-dependent early deposition can still become
imbalanced.

\paragraph{Cell-decomposition parallelism.}
Cell decomposition can be useful for a fully coupled Eulerian--Lagrangian
solver or for an expensive Eulerian flow solve.  It is less attractive here
because the carrier field is not solved and the particles are the dominant
work.  Cell decomposition also requires processor-boundary particle migration
and standard reconstruction workflows, which add overhead for this fixed-flow
one-way tracking use case.

\section{Conclusion}

For the impactor benchmark, \texttt{analyticPositionTracking} improves the
coarse-time-step behavior of the particle fate prediction.  It reduces the
large-step outlet/substrate bias and allows larger particle steps to be used
with less loss of accuracy.  The most sensitive size range in the present case
is the small-particle side, especially \(0.10\)--\(0.40\,\si{\micro\metre}\),
with an additional localized sensitivity around
\(0.85\,\si{\micro\metre}\) in some runs.

The comparison with \texttt{icoUncoupledKinematicParcelFoam} supports the
correctness of the implementation.  When
\texttt{analyticPositionTracking} is disabled, the custom solver reproduces the
legacy solver's fine-step patch fate counts.  The analytic-position branch then
improves the coarse-step behavior without changing the underlying particle
model.

For production calculations of this type, replicated-mesh particle sharding is
a strong parallelization strategy.  It directly divides the particle work and
avoids processor-patch particle migration.  The trade-off is increased memory
usage and sharded post-processing output.

\end{document}
